% Generated draft: requires mathematical and source review.
\documentclass[11pt]{article}
\usepackage[margin=25mm]{geometry}
\usepackage{fontspec}
\setmainfont{Latin Modern Roman}
\usepackage{amsmath,amssymb,mathtools,unicode-math}
\setmathfont{Latin Modern Math}
\usepackage{enumitem,needspace,xurl,hyperref}
\setlength{\parindent}{0pt}
\setlength{\parskip}{5pt}
\setlength{\emergencystretch}{3em}
\newcommand{\N}{\mathbb{N}}
\newcommand{\Z}{\mathbb{Z}}
\newcommand{\Q}{\mathbb{Q}}
\newcommand{\R}{\mathbb{R}}
\newcommand{\C}{\mathbb{C}}
\newcommand{\F}{\mathbb{F}}
\newcommand{\A}{\mathbb{A}}
\newcommand{\PP}{\mathbb{P}}
\newcommand{\E}{\mathbb{E}}
\newcommand{\eps}{\varepsilon}
\newcommand{\dd}{\,\mathrm{d}}
\begin{document}
\section*{UnsolvedMath ID: 30001078}
\textbf{Source problem code:} catalogue entry
\subsection{Definitions and mathematical statement}
The source asks two questions about random visibility: describe the boundary of the set $T$ of points visible from the origin in the Boolean model more precisely, and investigate maximal visibility in other random media, including bicolored random tessellations.

For orientation, a usual Boolean obstacle model has occupied set $\mathcal O=\bigcup_i(X_i+K_i)$ formed from a marked Poisson process. With the origin in the vacant set, a conventional visible set and directional visibility are
\[
T=\{x:[0,x]\cap\mathcal O=\varnothing\},\qquad
V(\theta)=\sup\{r\geq0:[0,r\theta]\cap\mathcal O=\varnothing\},
\]
for directions $\theta$ on the unit sphere in the intended dimension. Maximal visibility is then $V_{\max}=\sup_\theta V(\theta)$, possibly infinite. These displays are conditional definitions: the record does not give the intensity, grain distribution, dimension, conditioning on a vacant origin or endpoint convention. Allowing contact with obstacle boundaries can change whether visible segments include their endpoints.

The first request concerns the topological boundary $\partial T$, with a more detailed geometric or probabilistic description than the source currently provides. It specifies no particular dimension, regularity, distributional statistic or asymptotic scaling to be established. The second requests corresponding information on maximal visibility beyond the Boolean setting. For bicolored tessellations this requires a tessellation law, coloring law, which color obstructs sight and a sampling/conditioning convention; none is fixed.

The supplied assessment checked 2026-08-21 labels the questions open, refers to boundary and visibility questions for Poisson-line cells, and says existing visibility asymptotics do not give the requested general description. The main statement names a Boolean model, whereas this assessment names Poisson-line cells; their intended relationship is not explained. They are not silently treated as the same model. The corrected-verified note reports semantic and typographic cleanup without unsupported content additions.
\subsection{Short English statement}
Describe the geometry of the Boolean-model visible-set boundary and study maximal visibility in other random media, such as bicolored tessellations. The random-medium laws, visibility conventions and intended quantitative description are unspecified, and the assessment refers to a different Poisson-line setting.
\subsection{Sources}
\begin{itemize}
\item Local UnsolvedMath registry, permanent ID 30001078, source code catalogue entry; supplied complete record and corrected-verified cleanup note. Literature checked 2026-08-21: the assessment is supplied dated curation, not independently verified here.
\item New Perspectives in Stochastic Geometry (2009), Oberwolfach Reports, DOI 10.4171/owr/2008/47, pp. 2655--2702: \url{https://doi.org/10.4171/owr/2008/47}. Proposed-year metadata gives 2008.
\item Asymptotics of visibility: \url{https://citeseerx.ist.psu.edu/document?doi=4c4bdaaef250b3d96b41cca77c268189c859feb3&repid=rep1&type=pdf}.
\end{itemize}
\end{document}
